\chapter{The search region and the Sturm step}\label{chap:region}

The caller of \verb|find_points| may restrict the search to \emph{search intervals} for~$x/z$,
closed intervals whose ends are C \verb|double|s, and to a range of the denominators,
from~$b_{\mathrm{low}}$ to~$b_{\mathrm{high}}$. Chapter~\ref{chap:spec} specifies this search region
(Definition~\ref{def:regionSpec}); this chapter models what the C code does with it: the doubles,
the default interval and the refusal of bad intervals, the exact bounds of the numerators of each
denominator, and the Sturm step of \verb|sturm.c|, which narrows the search intervals to those on
which~$f$ may take a value $\ge 0$; and it proves, assuming Sturm's theorem
(Definition~\ref{def:sturmTheorem}), that the Sturm step keeps every point and that its intervals
fit the arrays that the C code gives it. The remainder chain of~$f$ and~$f'$, with which the Sturm
step begins, is Chapter~\ref{chap:setup}'s; the main theorems (Chapter~\ref{chap:theorems}) use
this chapter through the region that the search uses (Section~\ref{sec:searchRegion}).

\ratpoints{} does no floating-point arithmetic on the ends of its intervals (Chapter~\ref{chap:spec}).
So the model holds a double as its bits, with a function for each operation that the C code applies
to it, and computes everything else exactly; only the value of a double and well-formedness, the
test of \verb|RATPOINTS_BAD_ARGS|, read the bits.

\section{The doubles}\label{sec:doubles}

\begin{definition}[The doubles of the C code]\label{def:Double}
  \lean{Ratpoints.Double, Ratpoints.Double.lt, Ratpoints.Double.toRat, Ratpoints.Double.ofRatio, Ratpoints.Double.ldexp, Ratpoints.Double.nextUp}
  \leanok
  A C \verb|double| by its~$64$~bits of IEEE~754 binary64, in a structure with no equality of its
  own (equality of the bits is not \verb|==|); and the operations that \ratpoints{} applies to
  doubles, as gcc and glibc compute them with rounding to nearest, ties to even: the comparisons,
  false when either side is a NaN and otherwise the order of a key (the bits without the sign bit,
  negated when it is set); negation and \verb|fabs|; the exact value of a finite double; the double
  nearest to $\pm m\,2^k$ and to~$\pm n/d$; the conversions from and to \verb|long|, and to
  \verb|unsigned long|; and \verb|frexp|, \verb|ldexp| and \verb|nextafter| towards
  $\pm$\verb|HUGE_VAL|, which pass a signalling NaN on as it is, where glibc quiets it (\ratpoints{}
  passes them no NaN). These functions model the platform and hold no C range.
\end{definition}

Each operation is checked when the model is compiled, on the cases at the edges (zeros,
subnormals, the largest finite doubles, infinities, a NaN, ties, $2^{53}$ and~$2^{63}$ with their
neighbours), against the bits that a C program built with gcc and glibc prints; for the roundings
the proofs show what the theorems need, a faithful rounding.

\begin{lemma}[The values of the doubles]\label{lem:doubleValues}
  \uses{def:Double, def:regionSpec}
  \lean{Ratpoints.Double.lt_eq_true_iff, Ratpoints.Double.doubleValue_eq_toRat, Ratpoints.Double.ofDyadic_exact, Ratpoints.Double.ofInt64_exact, Ratpoints.Double.toRat_frexp_mul_zpow_eq_and_bounds, Ratpoints.Double.ldexp_exact, Ratpoints.Double.toInt_toInt64, Ratpoints.Double.toNat_toUInt64, Ratpoints.Double.nextUp_le_of_lt, Ratpoints.Double.lt_nextUp_and_nextDown_lt_ofRatio, Ratpoints.Double.isFinite_and_le_toRat_ofInt64_nextUp}
  \leanok
  Write~$v(x)$ for the value of the bits of a double~$x$ (Definition~\ref{def:regionSpec}).
  \begin{enumerate}
    \item On doubles that are not NaNs, \verb|<|, \verb|<=| and \verb|==| are the order and the
      equality of the values; the value of a finite double is its exact value.
    \item The double nearest to $\pm m\,2^k$ is exact for $0 < m < 2^{53}$, $k \ge -1074$ and
      $m\,2^k < 2^{1024}$, and \verb|(double)n| for $|n| \le 2^{53}$; for a finite $x \ne 0$,
      \verb|frexp| writes~$x$ as $g\,2^e$ with $1/2 \le |g| < 1$, and \verb|ldexp(x, n)| is
      $x\,2^n$ when $2^{-1022} \le |x|\,2^n < 2^{1023}$; and a double $0 \le x < 2^{63}$ converts
      to the \verb|long| $\lfloor x \rfloor$, one $0 \le x < 2^{64}$ to the \verb|unsigned long|
      $\lfloor x \rfloor$.
    \item The next double above~$x$, not a NaN or~$+\infty$, is at most every double above~$x$;
      dually below.
    \item The roundings of~$\pm n/d$ and of $\pm m\,2^k$ are never a NaN, and for $n, m \ne 0$ a
      finite result is a faithful rounding: the value lies strictly between the values of its two
      neighbours.
    \item For $1 \le H < 2^{63}$, the double~$h$ of~$H$, moved up to the next double when
      $h < 2^{63}$ and the \verb|long| of~$h$ is below~$H$, is finite and $H \le v(h)$.
  \end{enumerate}
\end{lemma}

\begin{proof}
  \leanok
  The value is strictly increasing in the key; the comparisons and the neighbours follow. A
  significand $M < 2^{53}$ with the exponent~$q$ of its last place has the value $\pm M\,2^q$:
  item~(2). For item~(4) the rounded significand lies within half a unit of the exact quotient, the
  neighbours one unit away, also across a change of the exponent; item~(5) follows.
\end{proof}

\section{The search intervals}\label{sec:searchIntervals}

The library normalizes the range of the denominators and the depth \verb|sturm| of the Sturm step
in place and restores those of the caller afterwards (\csource{find_points.c}{4542}{4545},
\csource{find_points.c}{4560}{4563}); the model returns the normalized options instead. The search
intervals and the Sturm step come after the reversal check (Chapter~\ref{chap:setup}).

\begin{definition}[The normalized options and the search intervals]\label{def:searchIntervals}
  \uses{def:Options, def:Double}
  \lean{Ratpoints.Options.normalize, Ratpoints.searchIntervals}
  \leanok
  For the options of a search (Definition~\ref{def:Options}) and the height bound $H \ge 1$, the
  \emph{normalized} options: the least denominator~$1$ for one below~$1$; the largest
  denominator~$H$ for one below~$1$ or above~$H$; the depth of the Sturm step capped at
  $\mathrm{DBL\_MANT\_DIG} + 2 = 55$, beyond which the ends of the intervals of the step would not
  all be doubles; and the switch of the reversal turned off, for the private bit
  \verb|RATPOINTS_NO_REVERSE_AUTO|, when the caller gives search intervals or the normalized range
  is not $1, \dots, H$ (\csource{find_points.c}{4659}{4661}, \csource{find_points.c}{4678}{4689}).
  The \emph{search intervals} of a search: for none given, the default interval $[-h, h]$, where~$h$
  is the double of~$H$, moved up to the next double when $h < 2^{63}$ and the \verb|long| of~$h$ is
  below~$H$; otherwise the caller's intervals, or none, for \verb|RATPOINTS_BAD_ARGS|, when an
  interval has not $l \le u$ or does not begin above the end of the one before it, compared as
  doubles, so that a NaN end is refused too (\csource{find_points.c}{4831}{4855}).
\end{definition}

The search turns the answer none into the refusal $\mathrm{badArgs}$; the test of a null pointer
(\csource{find_points.c}{4830}{4830}) is excluded as part of the interface of the library.

\begin{lemma}[The normalized options and the search intervals]\label{lem:searchIntervals}
  \uses{def:searchIntervals, def:regionSpec, def:Options}
  \lean{Ratpoints.Options.reverse_normalize, Ratpoints.Options.toNat_bHigh_normalize, Ratpoints.Options.sturm_normalize_le, Ratpoints.run_searchIntervals_eq_none_iff, Ratpoints.intervalsWF_run_searchIntervals, Ratpoints.exists_run_searchIntervals_nil}
  \leanok
  \begin{enumerate}
    \item The normalized switch of the reversal is the switch and the whole box
      (Definition~\ref{def:regionSpec}); the normalized largest denominator is
      $\mathrm{denominatorBound}(b_{\mathrm{high}}, H)$; a denominator $1 \le z \le H$ lies in the
      normalized range exactly when $b_{\mathrm{low}} \le z$ and ($b_{\mathrm{high}} < 1$ or
      $z \le b_{\mathrm{high}}$), as in the search region; the normalized depth is at most~$55$ and
      negative exactly when the depth is.
    \item The search intervals are none exactly when the caller's intervals are not well formed;
      they are the caller's intervals when these are well formed and not empty; and for
      $1 \le H < 2^{63}$ they are well formed and, when the caller gives none, hold~$a/b$ for every
      $b > 0$ and $|a| \le H$.
  \end{enumerate}
\end{lemma}

\begin{proof}
  \uses{lem:doubleValues}
  \leanok
  Item~(1) by cases on the clamps. Item~(2): the loop of the validation tests each index, and the
  chain of the specification, read index by index, is that test in the order of the values
  (Lemma~\ref{lem:doubleValues}); the default interval by item~(5) of that lemma.
\end{proof}

\section{The bounds of the numerators}\label{sec:numeratorBounds}

For a denominator~$b$ the numerators of an interval $[l, u]$ are $\lceil b l \rceil, \dots,
\lfloor b u \rfloor$ within the height bound. The C code computes them exactly from the doubles of
the ends, on \verb|long|s and in~$128$~bits where those do not suffice, and the model on the same
words; the bounds are made once for the run (\csource{find_points.c}{3869}{3870}).

\begin{definition}[The exact bounds of the search intervals]\label{def:exactBounds}
  \uses{def:Double, def:packing, def:bitArrays, def:reciprocal}
  \lean{Ratpoints.Sieve.exactEnd, Ratpoints.Sieve.exactMantissa, Ratpoints.Sieve.exactScale, Ratpoints.Sieve.endForm, Ratpoints.Sieve.exactBounds, Ratpoints.Sieve.intervalBits, Ratpoints.Sieve.intervalBitsExact, Ratpoints.Sieve.RangesFor}
  \leanok
  \begin{enumerate}
    \item An \emph{end}~$x$ (\csource{find_points.c}{3264}{3271},
      \csource{find_points.c}{3272}{3291}, \csource{find_points.c}{3292}{3299}): its sign, $x < 0$
      as in C, so that~$-0.0$ is not negative; for $|x| < 2^{63}$, $|x| = m\,2^{-s}$ with~$m$ odd
      or $s = 0$, except that~$s$ is capped at~$127$, where $m\,2^{-s} \ne |x|$ but $b\,m\,2^{-s}$
      and~$b|x|$ both lie strictly between~$0$ and~$1$ for $1 \le b < 2^{63}$; otherwise
      $m = 2^{63}$ and $s = 0$; and the rounding $r = 2^s - 1$ when the end wants~$b|x|$ rounded up
      (a lower end with $x \ge 0$, an upper end with $x < 0$), else $r = 0$. Its \emph{scaled} value
      at~$b$, $\min(2^{64} - 1, \lfloor (b m + r)/2^s \rfloor)$, is computed on two words
      (\csource{find_points.c}{3231}{3238}).
    \item The \emph{form} $(M, C, \mathit{sh})$ of an end (\csource{find_points.c}{3304}{3318}),
      for which $(b M + C) \gg \mathit{sh}$ fits a \verb|long| for every $b \le b_{\mathrm{high}}$:
      $(0, \mp H, 0)$ for an end at or beyond the height bound on its side ($x \le -H$ at a lower
      end, $x \ge H$ at an upper end); $(\pm m, C, s)$, with $C = 2^s - 1$ at a lower end and~$0$
      at an upper end, when $s < 64$ and $m \le \lfloor (2^{63} - 1 - C)/b_{\mathrm{high}}
      \rfloor$; and none otherwise.
    \item The \emph{bounds} of the search intervals (\csource{find_points.c}{3323}{3340}): the two
      ends of each interval, which is taken inline, with the forms of its ends, when
      $H \le 2^{63} - 65$ and both ends have a form, and otherwise gets the forms of an empty
      interval.
    \item The \emph{bits} of an interval at~$b$ for a class of the numerators
      (\csource{find_points.c}{3410}{3434}): $\mathit{lo} = \lceil b l \rceil$ and
      $\mathit{hi} = \lfloor b u \rfloor$ from the forms; the code~$2$ for $\mathit{lo} > H$; for
      $\mathit{hi} < -H$ the code~$1$ if the interval is taken inline, and otherwise what
      \verb|interval_bits_exact| computes from the scaled ends
      (\csource{find_points.c}{3381}{3406}); and otherwise the code~$0$ with the bits of the first
      and the last numerator of the class in $\max(\mathit{lo}, -H), \dots, \min(\mathit{hi}, H)$.
    \item The bounds of an interval \emph{represent} a range $\mathit{lo}, \dots, \mathit{hi}$
      at~$b$ when, for every class of a packing (Chapter~\ref{chap:sieve}: stride~$2^k$ with
      $k < \mathrm{numStrides}$ and $k \le 6$, residue $a_0 < 2^k$), their bits have the code~$2$
      only for a range above~$H$, the code~$1$ only for one below~$-H$, and otherwise the bits of
      the numerators of the class in $\max(\mathit{lo}, -H), \dots, \min(\mathit{hi}, H)$. The
      bounds of a run \emph{give} a list of ranges at~$b$ when there is one for each interval,
      represented by its bounds, each with $\mathit{lo} \le \mathit{hi} + 1$, increasing and
      strictly apart.
  \end{enumerate}
\end{definition}

This is what the loop of the sieve of one denominator over the intervals needs
(Chapter~\ref{chap:sieve}, \csource{find_points.c}{3471}{3488}): the code~$2$ ends the denominator,
the code~$1$ goes on to the next interval and the code~$0$ sieves the bits, so that the numerators
of a denominator come in increasing order.

\begin{lemma}[The bounds give the ranges of the numerators]\label{lem:numeratorRanges}
  \uses{def:exactBounds, def:regionSpec, def:Options}
  \lean{Ratpoints.Sieve.numeratorRanges, Ratpoints.Sieve.inRanges_numeratorRanges, Ratpoints.Sieve.le_add_one_of_mem_numeratorRanges, Ratpoints.Sieve.pairwise_numeratorRanges, Ratpoints.Sieve.lowNumerator_eq_and_highNumerator_eq_run_exactEnd, Ratpoints.Sieve.represents_intervalBoundsOf, Ratpoints.Sieve.rangesFor_run_exactBounds}
  \leanok
  The \emph{range of the numerators} of~$b$ in an interval $[l, u]$ over the values runs
  from~$\lceil b l \rceil$, clamped to $-H, \dots, H + 1$, to~$\lfloor b u \rfloor$, clamped to
  $-H - 1, \dots, H$, an infinite end going to the end of its clamp.
  \begin{enumerate}
    \item For $b > 0$ and $|a| \le H$, $a$ lies in one of the ranges of the intervals of~$D$
      exactly when~$a/b$ lies in one of the intervals. For intervals with $l \le u$ each range has
      $\mathit{lo} \le \mathit{hi} + 1$, and for search intervals (Definition~\ref{def:regionSpec})
      each range ends below the start of the next.
    \item For a double~$x$ that is not a NaN, $1 \le b < 2^{63}$ and $H < 2^{63}$, the end made
      of~$x$, at either side, gives the same clamped numerators as~$x$; and for $H < 2^{63}$ and
      $1 \le b \le b_{\mathrm{high}} \le H$ the bounds of an interval represent the range of the
      ends made of its doubles, on both paths of \verb|interval_bits|.
    \item So, for well-formed intervals, $H < 2^{63}$, $b_{\mathrm{high}} \le H$ and
      $0 < b \le b_{\mathrm{high}}$, the bounds of the run give at~$b$ the ranges of the values of
      the intervals (Definition~\ref{def:exactBounds}).
  \end{enumerate}
\end{lemma}

\begin{proof}
  \uses{lem:doubleValues, def:bitArrays, def:packing, def:reciprocal, lem:packing, lem:reciprocal}
  \leanok
  Item~(1): the ceiling and the floor of~$b x$ against an integer; the ranges are apart since
  $\lfloor b u_i \rfloor \le b u_i < b l_{i+1} \le \lceil b l_{i+1} \rceil$. Item~(2): the product
  on two words with its carry gives the scaled end exactly for $s \le 127$
  (\lname{Ratpoints.Sieve.toNat\_exactScale}); the mantissa is exact by
  Lemma~\ref{lem:doubleValues}, and an end capped at $s = 127$ or beyond~$2^{63}$ gives the clamped
  numerators of the double; a form $(b M + C) \gg s$ is $\lceil b m\,2^{-s} \rceil$ for
  $C = 2^s - 1$ and the floor for $C = 0$, a \verb|long| by the bound on~$m$. Item~(3): the loop of
  \verb|exact_bounds| is a map.
\end{proof}

\section{The Sturm step}\label{sec:sturmStep}

With $\mathrm{sturm} \ge 0$, \verb|_ratpoints_compute_sturm| narrows the search intervals before the
sieve. From the remainder chain $s_0 = f, s_1 = f', \dots, s_k$ (Chapter~\ref{chap:setup}) it
counts the sign changes at $\pm\infty$; when~$f$ has no real root and is negative it stops;
otherwise it bisects $[-\infty, +\infty]$ to the depth of the step, keeping the intervals on
which~$f$ may be $\ge 0$, and intersects them with the search intervals. The model keeps the
\verb|long| numerators and exponents of the ends, as \verb|Int64|, and the doubles of the C code.

\begin{definition}[The Sturm step]\label{def:computeSturm}
  \uses{def:Double, def:Options, def:remainderChain, def:reversal}
  \lean{Ratpoints.Sturm.computeSturm, Ratpoints.Sturm.iterate, Ratpoints.Sturm.EndsWF, Ratpoints.Sturm.intersect, Ratpoints.Sturm.meet, Ratpoints.sturmRegion}
  \leanok
  The \emph{bisection}, the nested function \verb|iterate| of gcc (\csource{sturm.c}{197}{208},
  \csource{sturm.c}{228}{233}, \csource{sturm.c}{261}{267}), works on intervals whose ends
  are~$n/2^e$, with \verb|long|s~$n$ and~$e$, or $\mp\infty$, written $\mp\mathrm{LONG\_MAX}$ with
  the exponent~$0$. It evaluates the rows of the chain at~$n/2^e$ by Horner's rule in integers,
  which computes~$2^{e \deg s}$ times the value of the row~$s$, and counts their sign changes, zeros
  skipped. An interval with the counts~$c_l$, $c_r$ at its ends and the signs~$\sigma_l$, $\sigma_r$
  of~$f$ just inside them is dropped when $c_l = c_r$ and $\sigma_l < 0$; it is kept when
  $c_l = c_r$ and $\sigma_l > 0$, or at the depth of the step, its ends converted to doubles,
  $\mathrm{ldexp}((\mathrm{double})\,n, -e)$, and appended to the intervals kept, or joined to the
  last of them when that ends at its lower end (\verb|==| of doubles); and otherwise it is split at
  its midpoint~$m$, the left half bisected first, both halves with the count and the sign of~$f$
  at~$m$, except that at a root~$m$ the left half gets the count plus~$1$ and
  $-\mathrm{sign}\, f'(m)$, the right half $\mathrm{sign}\, f'(m)$. The midpoint is~$0$ for
  $[-\infty, +\infty]$; $-1$ or~$2 n_r$ for $[-\infty, n_r]$ and~$1$ or~$2 n_l$ for
  $[n_l, +\infty]$, as the finite end is~$0$ or not; and otherwise half the sum of the ends over the
  larger denominator. An interval is \emph{well formed} at the depth~$\delta$ when it has one of the
  four shapes that the bisection meets (the whole line; an infinite interval whose finite end, of
  exponent~$0$, is~$0$ or on its side of~$0$; two finite ends, equal exponents being~$0$), its
  exponents lie between~$0$ and~$\delta$, its finite numerators over the larger exponent are at
  most~$2^{\delta}/4$ in absolute value from the depth~$1$ on, and its ends increase.

  The \emph{intersection} with the caller's intervals goes by two cursors: of the two current
  intervals, compared as doubles, the one that begins first is skipped when it ends before the other
  begins; otherwise their intersection is noted and the cursor of the one that ends first moves on
  (\csource{sturm.c}{289}{322}). The \emph{Sturm step} on~$f$ of degree $d \ge 1$ with $c_d \ne 0$
  (\csource{sturm.c}{86}{97}, \csource{sturm.c}{177}{190}, \csource{sturm.c}{269}{273}): none when
  the remainder chain is none; no interval when the counts at $\mp\infty$ agree, so that~$f$ has no
  real root, and $c_0 < 0$; otherwise the intervals that the bisection keeps from
  $[-\infty, +\infty]$, started with the counts and the signs of~$f$ at~$\mp\infty$, intersected
  with the caller's intervals. The \emph{Sturm step of the search}, on the polynomial that the
  reversal check hands on, at the normalized depth (\csource{find_points.c}{4879}{4892}): with the
  step off, the intervals as they are; the refusal $\mathrm{notSquarefree}$ when the chain is none;
  the intervals left, when there are some; and when there are none, none for an even~$d$ whose
  leading coefficient is not a square, so that the search returns at once, there being no point at
  infinity either, and otherwise no intervals with the flag \verb|sturm_empty|, by which the search
  delivers the points at infinity only.
\end{definition}

The proofs read~$f$ also as the real polynomial \lname{Ratpoints.Sturm.chainPoly} with the
coefficients $c_0, \dots, c_d$.

\begin{lemma}[What the Sturm step computes]\label{lem:sturmStep}
  \uses{def:computeSturm, def:regionSpec, def:Options, def:remainderChain, lem:squarefree_binaryFormPoly}
  \lean{Polynomial.forall₂_posAssoc_sturmSeq_of_getElem, Ratpoints.Sturm.IsChain.run_signsAt_eq, Ratpoints.Sturm.IsChain.run_countsAtInfinity_eq, Ratpoints.Sturm.iterate_wf, Ratpoints.Sturm.computeSturm_wf, Ratpoints.run_sturmRegion_eq_error_iff, Ratpoints.run_sturmRegion_subset}
  \leanok
  \begin{enumerate}
    \item Over an ordered field, a list of nonzero polynomials whose last entry is a unit and whose
      entries from the third on are \emph{positive multiples} ($c\,p$ with $c > 0$) of the negated
      remainders of the two before is, entry by entry, a list of positive multiples of the Sturm
      sequence of its first two entries; so the rows of a chain that ends with a constant are,
      over~$\mathbb{R}$, positive multiples of the Sturm sequence of~$f$ and~$f'$.
    \item Horner's rule gives the sign of the value of a row at~$n/2^e$ when $e \ge 0$ and~$e$
      times the size of the row is below~$2^{63}$. For a chain that ends with a constant, the
      counts of the step at $\pm\infty$ are those of the Sturm sequence of~$f$ and~$f'$; if also
      $d < 2^{57}$ and~$f$ is squarefree, the signs and the count of the step at~$n/2^e$,
      $0 \le e \le 55$, are the signs of~$f$ and~$f'$ and the count $V_f(n/2^e)$ of Sturm's
      theorem (Definition~\ref{def:sturmTheorem}).
    \item The whole line is well formed at the depth~$0$, and the halves of a well-formed interval
      at a depth below~$55$ are well formed at the next depth, without overflow of the \verb|long|s;
      the doubles of the ends have the values $n\,2^{-e}$, and~$\pm 2^{63}$ for
      $\pm\mathrm{LONG\_MAX}$. So the bisection to a depth of at most~$55$ keeps intervals with
      $l \le u$, increasing and strictly apart.
    \item The step is none exactly when the chain is none. The intersection of two well-formed
      lists holds every value that lies in both. At a depth of at most~$55$, for well-formed
      intervals of the caller, the step and the Sturm step of the search leave well-formed
      intervals inside the given ones, none when \verb|sturm_empty| is set; the Sturm step of the
      search refuses exactly when it is on and the chain is none.
  \end{enumerate}
\end{lemma}

\begin{proof}
  \uses{lem:doubleValues, def:reversal, lem:chain_squarefree}
  \leanok
  Item~(1): the chain has $A s_{j-2} = Q s_{j-1} - g s_j$ with integers $A, g \ge 1$
  (Chapter~\ref{chap:setup}); the rest by induction along the Sturm sequence, which commutes with
  the map from~$\QQ$ to~$\mathbb{R}$. Item~(2): positive multiples have the signs of the values and
  of the leading coefficients, and a squarefree~$f$ and~$f'$ never vanish together, which the count
  with the zeros skipped needs. Item~(3): the cases of the midpoint, whose numerators stay at
  most~$2^{53}$, so that the doubles are exact (Lemma~\ref{lem:doubleValues}). Item~(4): the loop
  of the intersection as a recursion on the two lists.
\end{proof}

\section{Sturm's theorem and the coverage}\label{sec:sturmCoverage}

The specification states Sturm's theorem as a proposition (Definition~\ref{def:sturmTheorem}), in
the form for an interval $(a, b]$ whose ends may be roots, which the bisection needs, as it
evaluates the chain at midpoints that may be roots of~$f$. The general facts behind the proofs
(signs near a point and near $\pm\infty$, sign variations of positive multiples, squarefree
polynomials over a perfect field) are stated at Mathlib's generality under \verb|Ratpoints/Mathlib|.

\begin{lemma}[What Sturm's theorem says about an interval]\label{lem:sturmInterval}
  \uses{def:sturmTheorem}
  \lean{Ratpoints.vLeft, Ratpoints.vRight, Ratpoints.sgnLeftEnd, Ratpoints.sgnRightEnd, Ratpoints.vLeft_eq_vRight_iff, Ratpoints.eval_neg_of_vLeft_eq_vRight, Ratpoints.exists_inside_eval_pos, Ratpoints.eval_neg_of_vBot_eq_vTop}
  \leanok
  For a real polynomial~$p$ let~$V_p(t)$ be the number of sign changes at~$t$ of the Sturm sequence
  of~$p$ and~$p'$, $V_p(+\infty)$ that of the leading coefficients of its entries, and
  $V_p(-\infty)$ that of the leading coefficients each times~$(-1)^{\deg q}$ for the entry~$q$. For
  an interval with ends $l < r$ in $\mathbb{R} \cup \{\pm\infty\}$, the \emph{count} at~$l$
  is~$V_p(l)$, and at~$r$ it is~$V_p(r)$, plus~$1$ at a root; the \emph{sign} at a finite~$l$ is
  that of~$p(l)$, or of~$p'(l)$ when $p(l) = 0$, at a finite~$r$ that of~$p(r)$, or of~$-p'(r)$ when
  $p(r) = 0$, and at $\pm\infty$ that of the leading coefficient of~$p$, times~$(-1)^{\deg p}$
  at~$-\infty$; for a squarefree~$p$ these are the signs of~$p$ just inside the ends. Assuming
  Sturm's theorem, for a squarefree~$p$:
  \begin{enumerate}
    \item the counts at the ends agree exactly when~$p$ has no root strictly between them;
    \item when they agree and~$p$ is negative just inside~$l$, $p$ is negative on $(l, r)$ and just
      inside~$r$;
    \item when they differ, or~$p$ is positive just inside~$l$, $p$ is positive somewhere in
      $(l, r)$;
    \item when $V_p(-\infty) = V_p(+\infty)$ and $p(0) < 0$, $p$ is negative everywhere.
  \end{enumerate}
\end{lemma}

\begin{proof}
  \leanok
  Sturm's theorem on $(l, r]$ counts the roots, a root at~$r$ giving the~$1$; at $\pm\infty$ the
  counts and the signs are those at points far enough out. Without a root, a sign just right of~$l$
  holds on $(l, r)$ by the intermediate value theorem; a root in $(l, r)$ gives a point of positive
  value, since~$p$ changes its sign at a simple root.
\end{proof}

\begin{lemma}[The Sturm step keeps every point]\label{lem:sturmCoverage}
  \uses{lem:sturmInterval, lem:sturmStep, def:computeSturm, def:sturmTheorem, def:binaryForm, def:Options, def:remainderChain, lem:squarefree_binaryFormPoly}
  \lean{Ratpoints.Sturm.IterateInvariant, Ratpoints.Sturm.covers_iterate, Ratpoints.Sturm.squarefree_chainPoly, Ratpoints.Sturm.covers_computeSturm, Ratpoints.run_sturmRegion_covers, Ratpoints.Sturm.size_run_iterate_le_of_chain_eq_some, Ratpoints.Sturm.size_run_computeSturm_le}
  \leanok
  Let~$f$ have degree $1 \le d < 2^{57}$ with $c_d \ne 0$, and the depth be at most~$55$. When the
  chain does not fail, $f$ is squarefree. Assuming Sturm's theorem:
  \begin{enumerate}
    \item The \emph{invariant} of the bisection is a well-formed interval whose counts and signs
      are those of~$f$ at its ends (Lemma~\ref{lem:sturmInterval}). The bisection of an interval
      with the invariant keeps what the intervals kept before it hold and every $x \in \QQ$ with
      $f(x) \ge 0$ between the values of the doubles of its ends, at a finite end only when~$f$ is
      positive just inside it.
    \item For well-formed intervals of the caller, the step leaves intervals that hold every
      $x \in \QQ$ of the caller's intervals with $|x| \le 2^{63}$ and $f(x) \ge 0$. So the Sturm
      step of the search, for~$d$ the degree of~$c$ and when it does not refuse, puts every~$a/b$
      of the given intervals with $b > 0$, $|a| \le 2^{63}$ and $F(a, b) \ge 0$ into an interval
      that it leaves, \verb|sturm_empty| clear; this needs Sturm's theorem only when the step is
      on.
    \item The bisection keeps at most $1 + \lfloor d/2 \rfloor$ intervals, the size of the array
      \verb|ivlocal| of the C code, and the step leaves at most~$\lfloor d/2 \rfloor$ intervals
      more than the caller gives, within the room for $1 + \lfloor d/2 \rfloor$ more that
      \verb|sturm.c| asks of the caller's array.
  \end{enumerate}
\end{lemma}

\begin{proof}
  \uses{lem:sturmInterval, lem:sturmStep, def:reversal, lem:chain_squarefree}
  \leanok
  Item~(1) by the induction principle of the bisection: a dropped interval holds no~$x$ with
  $f(x) \ge 0$ (Lemma~\ref{lem:sturmInterval}(2)); at a split the step computes the signs and the
  count of~$f$ at the midpoint (Lemma~\ref{lem:sturmStep}), so that both halves have the invariant
  and one of them keeps each~$x$ that the interval keeps. Item~(2): the whole line has the invariant
  and keeps every $|x| \le 2^{63}$; at the exit $f < 0$ (Lemma~\ref{lem:sturmInterval}(4)); and
  $F(a, b) = b^{d + (d \bmod 2)} f(a/b)$. Item~(3): each kept interval has a point with $f > 0$
  (Lemma~\ref{lem:sturmInterval}(3)), $f < 0$ somewhere between two of them, and~$f$ changes its
  sign at most~$d$ times; an intersection has fewer intervals than its two lists together.
\end{proof}

\section{The region that the search uses}\label{sec:searchRegion}

The main theorems are stated over the search region of the caller
(Definition~\ref{def:regionSpec}); their proofs pass through the region that the search uses, the
intervals that the Sturm step leaves.

\begin{definition}[The region that the search uses]\label{def:searchRegion}
  \uses{def:searchIntervals, def:computeSturm, def:regionSpec, def:binaryForm, def:sturmTheorem, def:Options, lem:sturmCoverage, lem:searchIntervals, lem:sturmStep, def:remainderChain, def:reversal}
  \lean{Ratpoints.searchRegion, Ratpoints.searchRegion_eq_of_run_sturmRegion_eq, Ratpoints.searchRegion_subset, Ratpoints.searchRegion_eq_none, Ratpoints.searchRegion_covers}
  \leanok
  A definition of the proofs, not of the model: for the polynomial~$c$ of degree~$d$ that the
  reversal check hands on, $H$ and the options, none when the search intervals are none or the
  Sturm step of the search refuses or returns at once; no interval when it sets
  \verb|sturm_empty|; and otherwise the values of the intervals that it leaves. For well-formed
  intervals of the caller:
  \begin{enumerate}
    \item every~$x$ in the region lies in the caller's intervals, or the caller gives none;
    \item if the chain does not fail when the step is on, the region is none only for an even~$d$
      whose leading coefficient is not a square;
    \item assuming Sturm's theorem when the step is on, for $0 < d < 2^{57}$ the degree of~$c$,
      $1 \le H < 2^{63}$ and a chain that does not fail when the step is on: every~$a/b$ with
      $b > 0$, $|a| \le H$ and $F(a, b) \ge 0$ in the caller's intervals, or anywhere when there
      are none, lies in the region, which is not none.
  \end{enumerate}
\end{definition}

Items~(1) and~(2) follow from Lemmas~\ref{lem:searchIntervals} and~\ref{lem:sturmStep}, item~(3)
from Lemma~\ref{lem:sturmCoverage} and the default interval. With the exact check on, and assuming
Sturm's theorem when the step is on, the depth of the Sturm step does not change the list of points
that a search returns; Theorem~\ref{thm:order} states this beside the order of the points.

\section{Deviations}\label{sec:regionDeviations}

Two deviations of this part bear on what the proofs cover (Chapter~\ref{chap:ccode} lists all of
them by kind). The model of \verb|exact_scale| forms the $128$-bit product, sum and shift on
two~$64$-bit words, for the \verb|unsigned __int128| of the arm that the default build compiles;
the other arm, for a build without a $128$-bit type (\csource{find_points.c}{3239}{3256}), is
excluded. And the \verb|long|s are \verb|Int64|s, where a wrap would stand for an overflow of the C
code, which the bounds on the forms and on the ends of the bisection exclude;
Lemma~\ref{lem:sturmCoverage}(3) bounds the arrays of the Sturm step.
